\section{Results}

\subsection{Validation and contribution of post-processing}

Ideal arithmetic checks yield total variation distances from theory of
\(1.38 \times 10^{-9}\) for \(N = 21\) and \(2.50 \times 10^{-13}\) for \(N = 35\),
establishing functional correctness, not performance under noise.

In UC1/UC2, the selected SVMs achieve F1 scores of 0.919/0.923 and ROC areas of
0.965/0.958, but do not reduce operational cost. TOP-4 averages 1.00 iterations in both
scenarios versus 1.37/1.50 for TOP-1 (\(p_{\mathrm{Holm}} = 0.0024/0.0015\)); Method 2
requires 1.50/1.37, does not significantly improve TOP-1, and is worse than its TOP-4
ablation (\(p_{\mathrm{Holm}} = 0.0033/0.0096\)). The benefit therefore comes from
multi-candidate search. Acting on already-acquired data, the filter can discard
histograms that TOP-4 would use successfully. ZNE also fails to improve the reference
scenario: two and three levels average 3,072 and 4,813 shots, versus 1,399 for TOP-1 and
1,024 for TOP-4. These findings concern the tested methods, instance, and noise model;
they do not establish the general ineffectiveness of machine learning or mitigation.

\subsection{Memory protection and hybrid decoding}

Within Monte Carlo uncertainty, the repetition code reproduces
\(p_{L} = 3p^{2} - 2p^{3}\). Steane's physical-to-logical error relation is approximately
quadratic (logarithmic slope 1.94), and \(p_{L}\) intersects \(p\) near 0.08. This is a
pseudo-threshold for data-qubit errors with ideal syndromes, not the operational threshold
of a fault-tolerant architecture. In the surface code, increasing distance reduces
logical error below threshold and increases it above threshold; the sub-threshold fit
yields about 0.86\% in the \(Z\) basis and 0.81\% in the \(X\) basis for the adopted
memory circuit, noise, and decoder.

Under nominal noise, the stand-alone network does not significantly outperform MWPM in
the initial fourteen configurations. With adjacent-qubit correlations poorly represented
by the matching graph, however, the hybrid reduces logical error. At \(d = 3\),
\(p = 0.003\), and crosstalk intensity 0.01, \(p_{L}\) falls from 0.05829 with MWPM to
0.04851 with the hybrid: a 16.8\% relative reduction (Table 5.9 of the thesis), significant
by paired McNemar test. With \(d\) rounds, the benefit is smaller at distance 5 and not
significant at distance 7. In the tested distance-5 configurations, BP+OSD outperforms the
MWPM-based hybrid; varying distance and rounds separately also reveals the role of network
input size. The learned contribution thus depends on noise structure, analytical
baseline, and data representation, not distance alone.

\subsection{Factorization success and the uniform reference}

Without noise, single-shot success is 74.89\% (95\% confidence interval
{[}74.59\%; 75.18\%{]}), consistent with 75\%. As \(p_{g}\) grows, the trend is
compatible with a non-increasing function; at \(p_{g} = 0.5\), success is 24.59\%, near
the 24.61\% uniform reference (Section~\ref{sec:implementation-en}). Correct
factorizations alone therefore do not reveal how much quantum information remains.

\subsection{Effectiveness and limitations of the approximate QFT}

\textbf{In the Shor pilot for} \(N = 15\)\textbf{, the selected AQFT raises single-shot
success from 47.92\% to 62.79\%} (+14.87 percentage points on independent test data;
95\% Newcombe interval {[}12.73; 16.99{]}), while two-qubit ECR gates fall from 362 to
279 and depth from 1,351 to 1,141. The benefit concerns factorization, not only isolated
phase estimation.

For \(N = 21\), truncating QFTs inside the arithmetic reduces ECR gates from 49,661 to
23,178, but ideal success falls from 46.67\% to 40.13\%; all three noisy-comparison
intervals include zero. With 128 test shots per variant and noise level, sensitivity to
small effects is limited, so no observed advantage does not prove ineffectiveness. All
nine phase-estimation differences favor the selected variant and six individual 95\%
intervals exclude zero; the largest increase is 12.55 percentage points with ten qubits
({[}9.77; 15.30{]}). These comparisons are exploratory, without simultaneous correction
of the nine intervals. The best tested truncation degree is two or three and need not
remain constant for arbitrarily large registers.
